\chapter{Nếu bài này là một món ăn}
\begin{tinhhuongbox}{Tình huống | Situation}
\songngu{Đầu tiết mà hỏi kiến thức ngay, lớp sợ sai và chọn im lặng.}{If you ask content questions immediately at the start, the class fears being wrong and chooses silence.}
\songngu{Một câu hỏi hình ảnh sẽ mở lời dễ hơn.}{An image-based question opens speech more easily.}
\end{tinhhuongbox}
\begin{noitrenlop}
\songngu{Nếu bài này là một món ăn, nó là món gì, và đừng nói mì gói cho an toàn quá nhé.}{If this lesson were a dish, what dish would it be, and please do not say instant noodles just to stay safe.}
\end{noitrenlop}
\begin{vihieuqua}
\songngu{Học sinh thấy câu hỏi vui, nhưng vẫn phải giải thích lựa chọn của mình.}{Students find the question fun, but they still have to explain their choice.}
\end{vihieuqua}
\begin{englishpocket}
\songngu{Cụm nên nhớ: \textit{If this lesson were...}}{Phrase to keep: \textit{If this lesson were...}}
\end{englishpocket}
\chapter{Phần nào dễ hiểu nhất trước}
\begin{tinhhuongbox}{Tình huống | Situation}
\songngu{Khi bài mới hơi khó, hỏi chỗ khó nhất trước làm cả lớp co lại.}{When a new lesson feels difficult, asking about the hardest part first makes the whole class shrink.}
\songngu{Hỏi từ chỗ dễ sẽ mở được tiếng nói trước.}{Starting from the easiest part opens up speech first.}
\end{tinhhuongbox}
\begin{noitrenlop}
\songngu{Phần nào trong bài này em thấy dễ hiểu nhất trước, để lớp mình có chút tự tin đã?}{Which part of this lesson feels easiest to understand first, so our class can gain a little confidence?}
\end{noitrenlop}
\begin{vihieuqua}
\songngu{Cảm giác làm được là chiếc chìa khóa mở miệng học sinh.}{The feeling of being able to do it is the key that opens student voices.}
\end{vihieuqua}
\begin{englishpocket}
\songngu{Cụm nên nhớ: \textit{Gain a little confidence.}}{Phrase to keep: \textit{Gain a little confidence.}}
\end{englishpocket}
\chapter{Nếu giải thích cho bạn chậm hơn mình}
\begin{tinhhuongbox}{Tình huống | Situation}
\songngu{Muốn kéo học sinh khỏi kiểu hiểu cho mình, hãy buộc các em nói cho người khác hiểu.}{If you want students to move beyond understanding for themselves, make them speak so someone else can understand.}
\songngu{Đó là lúc câu nói bắt đầu có chất học thật.}{That is when speaking starts becoming real learning.}
\end{tinhhuongbox}
\begin{noitrenlop}
\songngu{Nếu phải giải thích cho một bạn đang lag nhẹ, em sẽ nói câu nào trước?}{If you had to explain this to a classmate who is lagging a little, what sentence would you say first?}
\end{noitrenlop}
\begin{vihieuqua}
\songngu{Câu hỏi này vừa có chút hài, vừa ép học sinh diễn đạt bằng lời của mình.}{This question is slightly funny, and it also forces students to explain in their own words.}
\end{vihieuqua}
\begin{englishpocket}
\songngu{Cụm nên nhớ: \textit{What would you say first?}}{Phrase to keep: \textit{What would you say first?}}
\end{englishpocket}
\chapter{Lỗi nào dễ mắc nhất}
\begin{tinhhuongbox}{Tình huống | Situation}
\songngu{Học sinh rất thích đoán lỗi vì lỗi nghe giống chuyện người thật hơn là sách thật.}{Students like guessing mistakes because mistakes sound more like real people than textbooks do.}
\songngu{Đoán lỗi là một cách kéo chú ý rất nhanh.}{Guessing mistakes is a very quick way to pull attention in.}
\end{tinhhuongbox}
\begin{noitrenlop}
\songngu{Nếu bài này muốn bẫy lớp, theo em nó sẽ bẫy ở đâu?}{If this lesson wanted to trap the class, where do you think it would set the trap?}
\end{noitrenlop}
\begin{vihieuqua}
\songngu{Câu hỏi này biến sự đề phòng thành sự tỉnh táo.}{This question turns caution into alertness.}
\end{vihieuqua}
\begin{englishpocket}
\songngu{Cụm nên nhớ: \textit{Where would it set the trap?}}{Phrase to keep: \textit{Where would it set the trap?}}
\end{englishpocket}
\chapter{Một câu trả lời kỳ cục nhưng hợp lý}
\begin{tinhhuongbox}{Tình huống | Situation}
\songngu{Nhiều học sinh im không phải vì không nghĩ ra, mà vì sợ nghĩ ra điều khác người.}{Many students stay silent not because they have no idea, but because they fear having a different idea.}
\songngu{Cho phép sự kỳ cục có lý sẽ mở được nhiều miệng hơn.}{Allowing strange but sensible answers opens more mouths.}
\end{tinhhuongbox}
\begin{noitrenlop}
\songngu{Ai có một câu trả lời hơi kỳ cục nhưng vẫn hợp lý, mời lên tiếng để cứu sự nhàm chán cho cả lớp.}{Who has an answer that is a little weird but still makes sense? Speak up and save the whole class from boredom.}
\end{noitrenlop}
\begin{vihieuqua}
\songngu{Lớp cảm thấy được nới sợ sai, nên mức tham gia tăng lên ngay.}{The class feels that the fear of being wrong has been loosened, so participation rises immediately.}
\end{vihieuqua}
\begin{englishpocket}
\songngu{Cụm nên nhớ: \textit{A little weird but still makes sense.}}{Phrase to keep: \textit{A little weird but still makes sense.}}
\end{englishpocket}